% code_and_data.tex
\label{sec:code_and_data}

Every load-bearing claim is backed by a runnable script in the companion
repository, and the paper's derived equations are generated from those scripts
(via \texttt{sympy.latex}) into \texttt{paper/sections/generated/} and
\verb|\input| where they appear, so the displayed equations cannot drift from
what was computed.

\paragraph{Derivations (\texttt{src/utilities/}).}
\texttt{electric\_induced\_action.py} is the core: it re-derives the magnetic
anchor $Q_B=\{4,4,16\}$ from first principles, the electric block
$\varepsilon=P=\{1,4,4\}$, the adjugate theorem
$\mu^{-1}=\operatorname{adj}(\varepsilon)$ (for general hop vectors), the doubled-root
dispersion and null split (symbolic, plus a $2\times10^4$-direction numeric sweep),
and the $O_h$ domain average. \texttt{trlnT\_prescription.py} establishes the
$\mathrm{Tr}\ln T$ mode-filling prescription (physical band); \texttt{ward\_safe\_
diagnosis.py} diagnoses the momentum-space band-sum divergence.

\paragraph{Engine cross-check (\texttt{src/experiments/}).}
\texttt{exp\_01\_electric\_permittivity\_extraction.py} reads $\varepsilon=P$ off the
dcl-core engine's real \texttt{HopOperator.step} output (importing
\texttt{dcl\_core.core3d}), the temporal-sector mirror of dcl-core \texttt{exp\_04};
its \texttt{data/} outputs record the extracted tensor and the checks.

\paragraph{Notes (\texttt{notes/}).}
\texttt{electric\_block\_derivation.md} is the narrative derivation;
\texttt{speed\_anisotropy\_and\_isotropy\_restoration.md},
\texttt{trlnT\_prescription.md}, and \texttt{ward\_safe\_vacuum\_polarization.md}
cover the open items. \texttt{paper/sections/audit\_table.tex}
(Table~\ref{tab:audit}) is the authority mapping each claim to its script.
